\documentclass[12pt,a4paper]{article}
\usepackage[margin=2.5cm]{geometry}
\usepackage{parskip}
\usepackage{hyperref}
\usepackage{amsmath}
\usepackage{booktabs}
\usepackage{array}
\usepackage{longtable}
\usepackage{graphicx}
\usepackage{setspace}
\usepackage{xcolor}
\usepackage{caption}
\onehalfspacing

\newcolumntype{L}[1]{>{\raggedright\arraybackslash}p{#1}}
\newcolumntype{C}[1]{>{\centering\arraybackslash}p{#1}}
\newcolumntype{R}[1]{>{\raggedleft\arraybackslash}p{#1}}

\title{\textbf{AMR Similarity Network --- Junior Handoff}\\
\large Technical Reference for Marc Jacob Doria \& Gerardo Luis Fernando\\[0.4em]
\normalsize Network Science in Paris Project, Pitch 1}
\author{Clark Rodriguez (NetSci Lead)}
\date{May 2026}

\begin{document}
\maketitle
\tableofcontents
\clearpage

% ============================================================
\section{Overview and Pipeline}
% ============================================================

This document is your reference for the network science component of
the AMR--climate project.
All numbers here were extracted directly from the processed data files;
nothing is estimated or approximated.

\subsection*{What the pipeline does}

\begin{enumerate}
  \item \textbf{Step 1 (01\_extract\_hadex3.py)} --- Extracts annual
    country-level means of TXx (hottest-day temperature) from HadEX3
    NetCDF files for 27 EU/EEA countries over 2000--2018.
  \item \textbf{Step 2 (02\_preprocess.py)} --- Cleans and merges
    ECDC AMR surveillance data (10 pathogen--antibiotic combinations)
    with the climate data to produce a master dataset of 3\,926 rows.
  \item \textbf{Step 3 (03\_network.py)} --- Bins countries into LOW/HIGH
    temperature groups; builds a $27 \times 10$ resistance matrix;
    constructs a cosine-similarity network; computes centrality metrics.
  \item \textbf{Step 4 (04\_analysis.py)} --- Mann--Whitney U tests,
    centrality boxplots, network map, temperature-vs-resistance scatter.
  \item \textbf{Step 5 (05\_network\_final.py)} --- Publication-quality
    network figure.
  \item \textbf{Step 6 (06\_community.py)} --- Louvain community
    detection.
\end{enumerate}

\subsection*{Key processed files}

\begin{tabular}{@{}ll@{}}
\toprule
File & Contents \\
\midrule
\texttt{data/processed/master\_dataset.csv} & 3\,926 rows, 27 countries, 16 columns \\
\texttt{data/processed/country\_temp\_groups.csv} & 27 countries, temperature group \\
\texttt{data/processed/resistance\_matrix.csv} & $27 \times 12$ (iso3, name, 10 combos) \\
\texttt{data/processed/centrality\_scores.csv} & 27 countries, 4 centrality metrics \\
\texttt{data/processed/mannwhitney\_results.csv} & 5 metrics, MWU statistics \\
\texttt{data/processed/community\_assignments.csv} & 27 countries, community 0--3 \\
\texttt{outputs/networks/amr\_similarity\_network.gexf} & Full network for Gephi \\
\bottomrule
\end{tabular}

% ============================================================
\section{Temperature Groups (LOW vs HIGH)}
% ============================================================

Countries were split at the median annual TXx (hottest-day temperature)
of \textbf{33.66\,\textdegree{}C} across the 27-country sample.
Countries below the median are labelled LOW; at or above are labelled HIGH.
This gives 13 LOW and 14 HIGH countries.

\begin{longtable}{@{}llrr@{}}
\caption{Country temperature groups ordered by mean TXx.}\label{tab:tempgroups}\\
\toprule
ISO3 & Country & Mean TXx (\textdegree{}C) & Group \\
\midrule
\endfirsthead
\multicolumn{4}{l}{\small\textit{Continued from previous page}}\\
\toprule
ISO3 & Country & Mean TXx (\textdegree{}C) & Group \\
\midrule
\endhead
\midrule
\multicolumn{4}{r}{\small\textit{Continued on next page}}\\
\endfoot
\bottomrule
\endlastfoot
ISL & Iceland         & 19.56 & LOW  \\
NOR & Norway          & 24.90 & LOW  \\
IRL & Ireland         & 26.61 & LOW  \\
SWE & Sweden          & 27.98 & LOW  \\
FIN & Finland         & 28.04 & LOW  \\
GBR & United Kingdom  & 28.14 & LOW  \\
EST & Estonia         & 29.19 & LOW  \\
LVA & Latvia          & 29.67 & LOW  \\
DNK & Denmark         & 30.56 & LOW  \\
LTU & Lithuania       & 30.79 & LOW  \\
POL & Poland          & 32.90 & LOW  \\
NLD & Netherlands     & 33.10 & LOW  \\
DEU & Germany         & 33.55 & LOW  \\
\midrule
BEL & Belgium         & 33.66 & HIGH \\
CZE & Czechia         & 33.76 & HIGH \\
AUT & Austria         & 33.95 & HIGH \\
SVK & Slovakia        & 34.18 & HIGH \\
SVN & Slovenia        & 34.32 & HIGH \\
HUN & Hungary         & 34.44 & HIGH \\
FRA & France          & 34.72 & HIGH \\
HRV & Croatia         & 34.77 & HIGH \\
ITA & Italy           & 34.98 & HIGH \\
ROU & Romania         & 35.05 & HIGH \\
BGR & Bulgaria        & 36.30 & HIGH \\
GRC & Greece          & 36.74 & HIGH \\
ESP & Spain           & 37.59 & HIGH \\
PRT & Portugal        & 37.94 & HIGH \\
\end{longtable}

% ============================================================
\section{Network Construction}
% ============================================================

\subsection*{Resistance matrix}

A $27 \times 10$ matrix was constructed where rows are countries and
columns are the 10 pathogen--antibiotic combinations with sufficient
ECDC coverage.
Each cell contains the country's mean percentage resistance across all
available years (2000--2018).
Cells with $>$50\% missing data caused the row to be dropped (no
countries were dropped); remaining NaN cells were imputed with the
column mean.

\subsection*{Cosine similarity}

For each pair of countries, we computed the cosine similarity of their
10-dimensional resistance vectors:
\[
  \cos(\mathbf{u}, \mathbf{v}) = \frac{\mathbf{u} \cdot \mathbf{v}}{\|\mathbf{u}\|\,\|\mathbf{v}\|}
\]
A value of 1.0 means identical profile shape; a value of 0 means
orthogonal (completely different) profiles.

\subsection*{Threshold selection}

\begin{table}[h]
\centering
\caption{Threshold comparison. The threshold 0.96 was selected.}
\label{tab:threshold}
\begin{tabular}{@{}rrrrl@{}}
\toprule
Threshold & Edges & Density & Isolates & Connected \\
\midrule
0.90 & 268 & 0.762 & 0 & Yes \\
0.93 & 167 & 0.475 & 0 & Yes \\
0.95 &  96 & 0.273 & 0 & Yes \\
\textbf{0.96} & \textbf{66} & \textbf{0.188} & \textbf{0} & \textbf{Yes} \\
0.97 &  40 & 0.114 & 3 & No  \\
\bottomrule
\end{tabular}
\end{table}

The threshold 0.96 was the highest value that kept the graph fully
connected (no isolated nodes) while keeping edges below 80,
giving a moderately sparse network suitable for community detection.

% ============================================================
\section{Centrality Metrics}
% ============================================================

\subsection*{Definitions for a biology audience}

\begin{description}
  \item[Degree centrality] Fraction of all possible countries that a
    given country is directly connected to (shares a cosine similarity
    $>$ 0.96 with).
    High degree = many similar-profile neighbours.
  \item[Betweenness centrality] How often a country lies on the
    shortest path between two other countries.
    High betweenness = the country acts as a ``bridge'' between
    otherwise disconnected clusters. \textit{Note: due to a weight
    interpretation issue in the software, betweenness values here
    should be treated as relative ranks, not absolute measures.}
  \item[Closeness centrality] How short a country's average path
    length is to all other countries.
    High closeness = the country is centrally reachable.
  \item[Eigenvector centrality] Weighted measure of how well-connected
    a country's neighbours are.
    High eigenvector = the country is connected to other well-connected
    countries (similar to the PageRank concept).
\end{description}

\begin{longtable}{@{}llrrrr@{}}
\caption{Full centrality scores for all 27 countries.}\label{tab:centrality}\\
\toprule
ISO3 & Group & Degree & Betweenness & Closeness & Eigenvector \\
\midrule
\endfirsthead
\multicolumn{6}{l}{\small\textit{Continued from previous page}}\\
\toprule
ISO3 & Group & Degree & Betweenness & Closeness & Eigenvector \\
\midrule
\endhead
\midrule
\multicolumn{6}{r}{\small\textit{Continued on next page}}\\
\endfoot
\bottomrule
\endlastfoot
POL & LOW  & 0.231 & 0.532 & 0.353 & 0.001 \\
SVK & HIGH & 0.077 & 0.517 & 0.348 & 0.003 \\
CZE & HIGH & 0.077 & 0.508 & 0.335 & 0.018 \\
FRA & HIGH & 0.154 & 0.492 & 0.316 & 0.123 \\
DEU & LOW  & 0.385 & 0.342 & 0.285 & 0.412 \\
HUN & HIGH & 0.308 & 0.314 & 0.315 & 0.000 \\
HRV & HIGH & 0.192 & 0.129 & 0.292 & 0.000 \\
EST & LOW  & 0.115 & 0.077 & 0.250 & 0.000 \\
IRL & LOW  & 0.192 & 0.052 & 0.271 & 0.205 \\
ROU & HIGH & 0.154 & 0.049 & 0.254 & 0.000 \\
BGR & HIGH & 0.231 & 0.049 & 0.308 & 0.000 \\
ITA & HIGH & 0.231 & 0.046 & 0.261 & 0.000 \\
NLD & LOW  & 0.269 & 0.009 & 0.237 & 0.346 \\
PRT & HIGH & 0.115 & 0.009 & 0.260 & 0.000 \\
GBR & LOW  & 0.192 & 0.006 & 0.233 & 0.232 \\
AUT & HIGH & 0.346 & 0.000 & 0.241 & 0.398 \\
ESP & HIGH & 0.077 & 0.000 & 0.251 & 0.000 \\
DNK & LOW  & 0.231 & 0.000 & 0.235 & 0.316 \\
BEL & HIGH & 0.192 & 0.000 & 0.270 & 0.205 \\
FIN & LOW  & 0.231 & 0.000 & 0.235 & 0.315 \\
LTU & LOW  & 0.192 & 0.000 & 0.301 & 0.000 \\
ISL & LOW  & 0.038 & 0.000 & 0.203 & 0.000 \\
GRC & HIGH & 0.077 & 0.000 & 0.212 & 0.000 \\
LVA & LOW  & 0.231 & 0.000 & 0.308 & 0.000 \\
NOR & LOW  & 0.231 & 0.000 & 0.235 & 0.316 \\
SVN & HIGH & 0.077 & 0.000 & 0.247 & 0.000 \\
SWE & LOW  & 0.231 & 0.000 & 0.235 & 0.316 \\
\end{longtable}

% ============================================================
\section{Statistical Results (Mann--Whitney U Tests)}
% ============================================================

Mann--Whitney U tests compare whether the distribution of each metric
is shifted between the LOW and HIGH temperature groups.
The test is non-parametric (no normality assumption) and appropriate
for the small group sizes here (13 vs.\ 14 countries).

\begin{table}[h]
\centering
\caption{Mann--Whitney U test results (two-sided). Significant at $\alpha = 0.05$.}
\label{tab:mwu}
\begin{tabular}{@{}lrrrrrc@{}}
\toprule
Metric & $U$ & $p$-value & LOW median & HIGH median & $n$ (L/H) & Sig? \\
\midrule
Resistance rate (\%) & 26.0  & 0.0017 & 15.40 & 29.55 & 13/14 & \textbf{Yes} \\
Eigenvector centrality & 133.0 & 0.0440 &  0.232 &  0.000 & 13/14 & \textbf{Yes} \\
Degree centrality   & 124.5 & 0.1025 &  0.231 &  0.154 & 13/14 & No \\
Closeness centrality & 57.0 & 0.1040 &  0.237 &  0.265 & 13/14 & No \\
Betweenness centrality & 71.0 & 0.3218 &  0.000 &  0.048 & 13/14 & No \\
\bottomrule
\end{tabular}
\end{table}

\subsection*{Linear regression: temperature vs resistance}

Ordinary least squares regression of country mean TXx
against country mean resistance rate:
\[
  R^2 = 0.504, \quad \text{slope} = 1.68\%/\text{\textdegree{}C},
  \quad p < 0.001
\]

% ============================================================
\section{Louvain Community Detection}
% ============================================================

The Louvain algorithm partitioned the 27-node network into
\textbf{4 communities} with a modularity of $Q = 0.506$.
Modularity above 0.5 is generally considered strong community structure.

\begin{table}[h]
\centering
\caption{Community composition summary.}
\label{tab:comm_summary}
\begin{tabular}{@{}crrrl@{}}
\toprule
Community & $n$ & LOW & HIGH & Key countries \\
\midrule
0 & 11 & 8 & 3 & NOR, IRL, SWE, FIN, GBR, DNK, NLD, DEU, BEL, AUT, FRA \\
1 &  2 & 0 & 2 & CZE, SVK \\
2 &  8 & 4 & 4 & LVA, LTU, POL, HRV, ITA, ROU, BGR, GRC \\
3 &  6 & 2 & 4 & ISL, EST, SVN, HUN, ESP, PRT \\
\bottomrule
\end{tabular}
\end{table}

\begin{longtable}{@{}llrr@{}}
\caption{Full community assignments.}\label{tab:comm_full}\\
\toprule
ISO3 & Country & Community & Group \\
\midrule
\endfirsthead
\multicolumn{4}{l}{\small\textit{Continued from previous page}}\\
\toprule
ISO3 & Country & Community & Group \\
\midrule
\endhead
\midrule
\multicolumn{4}{r}{\small\textit{Continued on next page}}\\
\endfoot
\bottomrule
\endlastfoot
NOR & Norway          & 0 & LOW  \\
IRL & Ireland         & 0 & LOW  \\
SWE & Sweden          & 0 & LOW  \\
FIN & Finland         & 0 & LOW  \\
GBR & United Kingdom  & 0 & LOW  \\
DNK & Denmark         & 0 & LOW  \\
NLD & Netherlands     & 0 & LOW  \\
DEU & Germany         & 0 & LOW  \\
BEL & Belgium         & 0 & HIGH \\
AUT & Austria         & 0 & HIGH \\
FRA & France          & 0 & HIGH \\
\midrule
CZE & Czechia         & 1 & HIGH \\
SVK & Slovakia        & 1 & HIGH \\
\midrule
LVA & Latvia          & 2 & LOW  \\
LTU & Lithuania       & 2 & LOW  \\
POL & Poland          & 2 & LOW  \\
HRV & Croatia         & 2 & HIGH \\
ITA & Italy           & 2 & HIGH \\
ROU & Romania         & 2 & HIGH \\
BGR & Bulgaria        & 2 & HIGH \\
GRC & Greece          & 2 & HIGH \\
\midrule
ISL & Iceland         & 3 & LOW  \\
EST & Estonia         & 3 & LOW  \\
SVN & Slovenia        & 3 & HIGH \\
HUN & Hungary         & 3 & HIGH \\
ESP & Spain           & 3 & HIGH \\
PRT & Portugal        & 3 & HIGH \\
\end{longtable}

% ============================================================
\section{Ready-to-Paste Results Paragraphs}
% ============================================================

Use the paragraphs below verbatim (or lightly edited) in any written
deliverable.
Check that the framing fits your section structure before pasting.

\subsection*{Primary finding: temperature and AMR resistance rate}

Among the 27 EU/EEA countries included in this analysis,
HIGH-temperature countries (those with mean annual TXx $\geq$ 33.66\,\textdegree{}C,
$n = 14$) exhibited significantly higher mean antimicrobial resistance rates
than LOW-temperature countries ($n = 13$)
(Mann--Whitney $U = 26.0$, $p = 0.0017$;
LOW median = 15.4\%, HIGH median = 29.5\%).
Linear regression of country-level mean TXx against mean resistance rate
confirmed a positive association ($R^2 = 0.504$, slope = 1.68\%/\textdegree{}C,
$p < 0.001$), indicating that approximately half of the cross-country variance
in resistance rates is associated with hottest-day temperature.

\subsection*{Secondary finding: eigenvector centrality}

Within the AMR cosine-similarity network (threshold $> 0.96$, 27 nodes,
66 edges, density = 0.188), LOW-temperature countries had significantly
higher eigenvector centrality than HIGH-temperature countries
(Mann--Whitney $U = 133.0$, $p = 0.044$;
LOW median = 0.232, HIGH median $\approx$ 0.000).
Eigenvector centrality quantifies the extent to which a node is connected
to other highly connected nodes; the result indicates that LOW-temperature
countries form a structurally cohesive hub within the network,
driven by the tightly interconnected western European cluster
(Germany, Austria, Netherlands, Denmark, Norway, Sweden, Finland,
and Ireland).
Degree, betweenness, and closeness centrality did not differ
significantly between temperature groups
($p = 0.103$, $0.322$, and $0.104$, respectively).

\subsection*{Community structure}

Louvain community detection partitioned the network into four
communities (modularity $Q = 0.506$).
Community 0 (11 countries) was dominated by LOW-temperature,
northern and western European nations, consistent with the eigenvector
centrality finding.
Community 1 consisted exclusively of Czechia and Slovakia---two
geographically adjacent HIGH-temperature countries with nearly identical
resistance profiles.
Communities 2 and 3 were climatically mixed, suggesting that additional
factors beyond temperature (antibiotic stewardship policies, healthcare
infrastructure, cross-border patient transfer) contribute to the
resistance profile clustering observed in eastern and southern Europe.

% ============================================================
\section{Figure List}
% ============================================================

All figures are saved in \texttt{outputs/figures/}:

\begin{table}[h]
\centering
\caption{Generated figures.}
\begin{tabular}{@{}ll@{}}
\toprule
File & Description \\
\midrule
\texttt{amr\_network\_final.png} & Publication-quality network coloured by temp group \\
\texttt{amr\_network\_map.png}   & Network map with betweenness-scaled nodes \\
\texttt{centrality\_boxplots.png} & $2 \times 2$ boxplots with MWU annotations \\
\texttt{temp\_vs\_resistance.png} & TXx vs mean resistance scatter + OLS line \\
\bottomrule
\end{tabular}
\end{table}

The GEXF file \texttt{outputs/networks/amr\_similarity\_network.gexf}
can be opened directly in Gephi for interactive exploration.
Node attributes included: \texttt{country\_name}, \texttt{temp\_group},
\texttt{mean\_TXx}, \texttt{degree}, \texttt{betweenness},
\texttt{closeness}, \texttt{eigenvector}, \texttt{community}.

% ============================================================
\section{Discussion Points}
% ============================================================

The following bullet points are intended to help you frame the results
in a biological/clinical context during discussions or presentations:

\begin{itemize}
  \item The positive association between temperature and resistance rate is
    ecologically plausible: higher ambient temperatures accelerate
    horizontal gene transfer between bacteria, increase community antibiotic
    use (driven by higher infectious disease incidence in warmer seasons),
    and are associated with longer outdoor survival of environmental
    resistant organisms.

  \item The 14-percentage-point gap in median resistance rates between
    LOW and HIGH groups (15.4\% vs.\ 29.5\%) is clinically meaningful.
    For a pathogen with a 20\% background resistance rate, this gap
    shifts the probability that an empirically chosen antibiotic will
    fail from roughly 1-in-5 to almost 1-in-3.

  \item The eigenvector centrality finding is structurally interesting:
    LOW-temperature countries cluster together in the similarity network,
    implying that cold-climate EU countries not only have lower resistance
    rates overall but also have more similar resistance profiles to each
    other.
    This may reflect shared antibiotic stewardship governance (Scandinavian
    countries have historically led EU antibiotic consumption reduction),
    or it may reflect a common selective environment in cooler climates.

  \item The communities identified by Louvain do not map cleanly onto
    temperature groups, particularly in eastern and southern Europe.
    This heterogeneity is consistent with the known role of healthcare
    system capacity and hospital infection control practices in shaping
    local resistance patterns---factors that vary substantially across
    the Balkans, for example, despite similar temperatures.

  \item The $R^2 = 0.504$ from the simple linear regression leaves
    roughly half the variance unexplained.
    Antibiotic consumption per capita (DDD/1000 inhabitants/day),
    which was not included in the network analysis, is a strong
    independent predictor of resistance rates in ECDC surveillance
    data and would be the most important covariate to add in a
    follow-up multivariate model.

  \item The 10 pathogen--antibiotic combinations retained in the
    resistance matrix represent surveillance-grade ECDC data
    (hospital-acquired isolates only), which is advantageous for
    comparability but does not capture community-acquired resistance
    or resistance in primary care settings.
\end{itemize}

% ============================================================
\section{Reproducing the Analysis}
% ============================================================

\begin{enumerate}
  \item Install dependencies: \texttt{pip install -r requirements.txt}
  \item Place HadEX3 \texttt{.nc.gz} files in \texttt{data/raw/hadex3/}
    and ECDC CSVs in \texttt{data/raw/ecdc/}.
  \item Run scripts in order:
    \begin{verbatim}
python scripts/01_extract_hadex3.py
python scripts/02_preprocess.py
python scripts/03_network.py
python scripts/04_analysis.py
python scripts/05_network_final.py
python scripts/06_community.py
    \end{verbatim}
  \item All outputs land in \texttt{data/processed/} and
    \texttt{outputs/}.
\end{enumerate}

Raw data files (\texttt{data/raw/}, \texttt{*.nc}, \texttt{*.nc.gz})
and processed CSVs are excluded from the repository by \texttt{.gitignore}.
The GEXF network file and pipeline scripts are tracked.

\end{document}
